\section{Síntesis y ampliación}

Esta sección condensa el EP110 en unas cuantas conclusiones. Primera: la definición mínima de un Agent es percepción y acción, no conversación. Segunda: Kimi K2, ChatGPT Agent, Qwen3-Coder y Manus muestran, respectivamente, las capas de modelo, producto, código/RL e ingeniería de contexto de la pila de un Agent. Tercera: el núcleo del Agent training pasa de input-output a trajectory, environment, task-reward y safety. Cuarta: lo decisivo en un Agent de producción es la ingeniería de sistemas: KV cache, enmascaramiento de herramientas, memoria en el sistema de archivos, recuperación ante errores y control por parte del usuario.

\begin{importantbox}{Ubicar el EP110 en la serie de IA de Zhang Xiaojun (张小珺)}
El EP115 presenta la teoría de la segunda mitad de los Agent; el EP113 aborda Kimi K2 desde la perspectiva de la empresa que desarrolla el modelo; el EP116 trata el Agentic Model empresarial. El EP110, en cambio, reúne varios informes técnicos recientes y muestra la cadena completa de un Agent, desde el artículo y el modelo hasta el producto y su puesta en práctica de ingeniería.
\end{importantbox}

\begin{knowledgebox}{Relación con los episodios anteriores y posteriores}
\begin{tabularx}{\textwidth}{p{0.18\textwidth}p{0.34\textwidth}X}
\toprule
Episodio & Tema & Conexión con el EP110 \\
\midrule
EP115 & Teoría de la segunda mitad de los Agent & Ofrece el marco teórico de reward, environment e interface. \\
EP113 & Kimi K2 y Agentic LLM & Muestra cómo una empresa de modelos entiende K2 y el ecosistema de código abierto. \\
EP116 & Agentic Model empresarial & Ofrece la perspectiva de despliegue de un Agent con datos privados en el ámbito ToB. \\
EP139 & Historia técnica de los Agent & Traza la evolución histórica de los Agent, desde los primeros sistemas hasta los LLM Agent. \\
\bottomrule
\end{tabularx}
\end{knowledgebox}

\subsection{Conclusiones clave}

Las secciones anteriores ya explicaron el Agent desde la definición, el entrenamiento, el producto, el entorno de código y la ingeniería de contexto. Esta sección condensa ese contenido en unos cuantos juicios de ingeniería, para facilitar su comparación posterior con el EP115, el EP113 y el EP116.

\begin{enumerate}
\item Agent = percepción + acción + retroalimentación; no es una capacidad aislada de «saber invocar herramientas».
\item Lo central de Kimi K2 es el modelo MoE abierto, MuonClip, las trayectorias sintéticas de uso de herramientas y el joint RL.
\item Lo central de ChatGPT Agent es el sistema de producto unificado y el control por parte del usuario.
\item Lo central de Qwen3-Coder es el entorno de código, el Code RL, el Long-horizon RL y la cadena de herramientas.
\item Lo central de Manus es la ingeniería de contexto, en particular el KV cache, el enmascaramiento de herramientas y la memoria en el sistema de archivos.
\end{enumerate}

\subsection{Preguntas abiertas}

Estas preguntas son los aspectos que siguen sin resolverse una vez que el Agent pasa de los informes al producto. De ellas dependen el próximo paradigma de entrenamiento, la arquitectura de producto y la dirección de la inversión en infraestructura.

\begin{enumerate}
\item En el entrenamiento de un Agent, ¿en qué proporción deben combinarse las trayectorias sintéticas y la retroalimentación de entornos reales?
\item ¿Pueden las Rubrics as reward sostener una self-improvement confiable en tareas abiertas?
\item ¿En qué forma convergerán finalmente las dos rutas, la in-context y la end-to-end?
\item En un Agent de producción, ¿el principal cuello de botella de costo vendrá de las llamadas al modelo, de la ejecución de herramientas o de la longitud del contexto?
\end{enumerate}

\subsection{Lecturas adicionales}

\begin{itemize}
\item Si le interesa el marco teórico de los Agent, puede compararse con la entrevista a Yao Shunyu (姚顺雨) del EP115.
\item Si le interesa la perspectiva de Kimi K2 como empresa de modelos, puede compararse con la entrevista a Yang Zhilin (杨植麟) del EP113.
\item Si le interesa el Agentic Model empresarial, puede compararse con la entrevista a Wu Minghui (吴明辉) del EP116.
\end{itemize}

\end{document}
